% =========================================
% CHAPTER 6: EXPERIMENTAL EVALUATION
% =========================================
\chapter{Experimental Results and Performance Evaluation}
\label{ch:evaluation}

This chapter presents the experimental methodology and results of the comparative evaluation. We describe the test environment and scenarios, and then analyze the results along four dimensions: throughput, batch efficiency, finality latency, and computational overhead. Each subsection includes a discussion of the observed trade-offs between the two architectures.

\section{Experimental Setup}

\subsection{Test Environment}

All experiments were conducted in a containerized environment with the following configuration:

\begin{itemize}
    \item \textbf{Hardware}: 8GB RAM, 4-core processor
    \item \textbf{OS}: Linux/Docker container environment
    \item \textbf{L1 Node}: Hardhat local node (production-like behavior)
    \item \textbf{Sequencers}: Synchronized with 5-second grace period
    \item \textbf{Duration}: Continuous 5-minute test scenarios
\end{itemize}

\subsection{Test Scenarios}

Three scenarios were defined to evaluate the framework under different load conditions:

\begin{enumerate}
    \item \textbf{Scenario~A (Normal Load)}: 150~transactions are submitted to each sequencer, representing a moderate workload intended to exercise the basic batching and submission pipeline.
    \item \textbf{Scenario~B (High Load)}: 500~transactions are submitted to stress-test the batching mechanism and measure peak throughput.
    \item \textbf{Scenario~C (Batch Analysis)}: 250~transactions are submitted to observe batching patterns and L1 submission frequency.
\end{enumerate}

Each scenario was run for a continuous 5-minute window, with both sequencers operating in parallel against the same L1 node.

\section{Throughput Analysis}

\subsection{Transactions Per Second}

Table~\ref{tab:tps-results} presents the measured peak TPS for each scenario.

\begin{table}[h]
    \centering
    \begin{tabular}{l|c|c|c}
        \hline
        \textbf{Scenario} & \textbf{ZK Rollup} & \textbf{Optimistic} & \textbf{Difference} \\
        \hline
        Normal (150~tx) & $\sim$25~TPS & $\sim$18~TPS & +39\% ZK \\
        \hline
        High Load (500~tx) & $\sim$50~TPS & $\sim$35~TPS & +43\% ZK \\
        \hline
        Batch (250~tx) & $\sim$40~TPS & $\sim$28~TPS & +43\% ZK \\
        \hline
    \end{tabular}
    \caption{Peak throughput comparison (transactions per second)}
    \label{tab:tps-results}
\end{table}

Across all three scenarios, the ZK~Rollup achieves 39--43\% higher peak throughput than the Optimistic~Rollup. This advantage arises from the ZK~sequencer's configuration:
\begin{itemize}
    \item a smaller batch size (100~transactions versus 200) triggers submissions more frequently;
    \item a shorter batching interval (5~seconds versus 8~seconds) reduces idle time between batches.
\end{itemize}

As a result, the ZK~sequencer processes transactions into L1-committed batches at a faster rate. However, this comes at the cost of a higher number of L1 submissions, which we analyze in the next subsection.

\subsection{Batch Efficiency}

While the ZK~Rollup achieves higher throughput, the Optimistic~Rollup makes more efficient use of each L1 submission. Table~\ref{tab:batch-efficiency} summarizes the batch efficiency metrics measured during a 400-transaction test.

\begin{table}[h]
    \centering
    \begin{tabular}{l|c|c|c}
        \hline
        \textbf{Metric} & \textbf{ZK Rollup} & \textbf{Optimistic} & \textbf{Trade-off} \\
        \hline
        Batch size & 100~txs & 200~txs & 2$\times$ difference \\
        \hline
        Batches created (400~tx) & 7 & 4 & 43\% fewer for Optimistic \\
        \hline
        L1 submissions & 7 & 4 & Higher L1 cost for ZK \\
        \hline
        L1 gas per transaction & Higher & Lower & Optimistic cheaper \\
        \hline
    \end{tabular}
    \caption{Batch efficiency comparison}
    \label{tab:batch-efficiency}
\end{table}

We define batch efficiency as the ratio of transactions processed to L1 submissions:
\begin{equation}
\text{Batch Efficiency} = \frac{\text{Transactions Batched}}{\text{L1 Submissions}}
\end{equation}

Under this metric, the ZK~Rollup achieves an efficiency of $400 / 7 \approx 57$~transactions per L1 submission, while the Optimistic~Rollup achieves $400 / 4 = 100$~transactions per L1 submission. This means that the Optimistic~Rollup reduces L1 interaction by approximately 43\% compared to the ZK~Rollup, which translates directly into lower per-transaction gas costs on L1.

This result illustrates a fundamental trade-off: the ZK~Rollup prioritizes low-latency finality (via frequent, smaller batches), while the Optimistic~Rollup prioritizes cost efficiency (via fewer, larger batches).

\section{Finality Analysis}

\subsection{L1 Inclusion Time}
Table~\ref{tab:l1-latency} shows the L1 inclusion latency---the time elapsed between the sequencer submitting a batch transaction and that transaction being mined into an L1 block---measured at several percentiles.

\begin{table}[h]
    \centering
    \begin{tabular}{l|c|c|c}
        \hline
        \textbf{Latency} & \textbf{ZK Rollup} & \textbf{Optimistic} & \textbf{L1 Baseline} \\
        \hline
        P50 (median) & 4.2s & 4.3s & 4.0s \\
        \hline
        P95 & 4.75s & 4.75s & 4.7s \\
        \hline
        P99 & 5.2s & 5.1s & 5.0s \\
        \hline
        Maximum & 8.3s & 8.1s & 8.0s \\
        \hline
    \end{tabular}
    \caption{L1 inclusion latency in seconds}
    \label{tab:l1-latency}
\end{table}

Both rollups achieve L1 inclusion in approximately 4.75~seconds at the P95 level, which is nearly identical to the L1 baseline (a plain Ethereum transaction on the same Hardhat node). This confirms that \textbf{L1 block production time is the dominant factor} in inclusion latency, and neither sequencer introduces significant additional delay at this stage.

\subsection{True Finality (Architecture-Specific)}

L1 inclusion is a necessary but not sufficient condition for finality. Each rollup architecture imposes additional requirements before a transaction can be considered truly final. Table~\ref{tab:true-finality} breaks down the complete finality path for each architecture.

\begin{table}[h]
    \centering
    \begin{tabular}{l|c|c|c}
        \hline
        \textbf{Finality Stage} & \textbf{ZK Rollup} & \textbf{Optimistic} & \textbf{Difference} \\
        \hline
        L1 inclusion & 4.75s & 4.75s & Identical \\
        \hline
        Proof verification & $\sim$15s & --- & ZK-specific \\
        \hline
        Challenge period & --- & 10~blocks ($\sim$2~min) & Optimistic-specific \\
        \hline
        \textbf{True finality} & \textbf{$\sim$15s} & \textbf{$\sim$2~min} & \textbf{$\sim$8$\times$ faster for ZK} \\
        \hline
    \end{tabular}
    \caption{True finality comparison in the demo environment}
    \label{tab:true-finality}
\end{table}

This distinction is critical. In our demo environment, the ZK~Rollup achieves true finality approximately 8 times faster than the Optimistic~Rollup. In a production Ethereum deployment, the gap widens dramatically:
\begin{itemize}
    \item \textbf{ZK~Rollup}: L1 inclusion ($\sim$12s) + proof verification ($\sim$2s) + confirmations ($\sim$1s) $\approx$ \textbf{15~seconds}.
    \item \textbf{Optimistic~Rollup}: L1 inclusion ($\sim$12s) + 7-day challenge period $\approx$ \textbf{604,800~seconds}.
\end{itemize}

This represents a difference of roughly 672$\times$ in finality time, which has direct implications for applications that depend on fast settlement (such as decentralized exchanges and cross-chain bridges).

\section{Computational Overhead: Proof Generation}

The ZK~Rollup incurs a computational cost that the Optimistic~Rollup avoids: generating a zero-knowledge proof for each batch. Table~\ref{tab:proof-latency} reports the measured proof generation latency.

\begin{table}[h]
    \centering
    \begin{tabular}{l|c|c}
        \hline
        \textbf{Statistic} & \textbf{Time (ms)} & \textbf{Interpretation} \\
        \hline
        Minimum & 45 & Low overhead \\
        \hline
        P50 (median) & 78 & Typical per-batch cost \\
        \hline
        P95 & 95 & Occasional higher latency \\
        \hline
        P99 & 120 & Rare spike \\
        \hline
        Average & 82 & Mean per-batch cost \\
        \hline
    \end{tabular}
    \caption{ZK proof generation latency distribution}
    \label{tab:proof-latency}
\end{table}

At the P95 level, proof generation takes 95~ms. Relative to the 5-second batching interval, this represents a modest overhead:

\begin{equation}
\text{Batch Overhead} = \frac{\text{Proof Time}}{\text{Batch Interval}} = \frac{95 \text{ ms}}{5000 \text{ ms}} = 1.9\%
\end{equation}

This indicates that, in our test environment, proof generation consumes less than 2\% of the available time per batch cycle. While production ZK proof systems operating over more complex circuits will incur significantly higher costs, this result demonstrates that proof generation does not represent a fundamental bottleneck in the architecture---particularly as proof systems continue to improve in efficiency.

\section{Security Model Comparison}

In addition to performance metrics, the evaluation also compared the security behavior of both architectures when confronted with an invalid state proposal. Table~\ref{tab:security-models} contrasts the key security properties.

\begin{table}[h]
    \centering
    \begin{tabular}{l|p{4cm}|p{4cm}}
        \hline
        \textbf{Property} & \textbf{ZK Rollup} & \textbf{Optimistic Rollup} \\
        \hline
        Security assumption & 1-of-N honest (one prover suffices) & At least one honest challenger required \\
        \hline
        Proof type & Cryptographic validity proof & Economic fraud proof \\
        \hline
        Invalid state halted & Immediately upon proof rejection & After fraud proof submission \\
        \hline
        Cost model & Computational (proof generation) & Economic (challenge period and bonds) \\
        \hline
        Censorship risk & Low (anyone can generate a proof) & Medium (requires active challengers) \\
        \hline
    \end{tabular}
    \caption{Security model comparison}
    \label{tab:security-models}
\end{table}

\subsection{Behavior Under Invalid State Proposals}

To illustrate the practical difference, consider a scenario in which the sequencer proposes an invalid state transition.

In the \textbf{ZK~Rollup}, the response proceeds as follows:
\begin{enumerate}
    \item The sequencer submits a proof together with the invalid state root.
    \item The L1 verifier contract checks the proof cryptographically.
    \item Since the proof does not match the claimed state root, the submission is rejected.
    \item The system continues from the last accepted valid state.
    \item Detection time: approximately 15~seconds (the next proof cycle).
\end{enumerate}

In the \textbf{Optimistic~Rollup}, the response differs significantly:
\begin{enumerate}
    \item The sequencer submits only the state root (no proof is attached).
    \item The state root is accepted as provisional.
    \item Any observer may submit a fraud proof if they detect that the state is invalid.
    \item The L1 contract re-executes the disputed transaction to adjudicate the claim.
    \item If fraud is confirmed, the invalid state is reverted and the sequencer is penalized.
    \item Time to resolution: the full challenge period (7~days in production).
\end{enumerate}

In summary, the ZK~Rollup detects invalidity through a cryptographic mechanism (fast but computationally expensive), while the Optimistic~Rollup relies on an economic mechanism (slower but cheaper in the absence of fraud).
